\chapter{Governing equations and constitutive relations}
\label{ch:equations}

This chapter states the field equations of the coupled model, in each case with the conservation
principle or constitutive theory from which the equation follows, the assumptions of
Chapter~\ref{ch:model} that make the stated form applicable, and the literature source of the
general relation. Symbol definitions and units are collected in Chapter~\ref{ch:model} and
Appendix~\ref{app:symbols}; the boundary and interface conditions are assembled separately in
Chapter~\ref{ch:interface}.

\section{Mass conservation in the fluid}
For an incompressible continuum, conservation of mass reduces to the divergence-free condition on
the velocity field \cite{bird1960transport},
\begin{equation}
\nabla\cdot\bu=0 \qquad \text{in } \Om_f .
\label{eq:continuity}
\end{equation}
Under Assumption~A2 the density is constant and \eqref{eq:continuity} holds identically. The domain
$\Om_f$ is the \emph{deformed} steady domain defined in Section~\ref{sec:ale}; the deformation
enters the fluid problem only through the domain geometry.

\section{Momentum conservation in the fluid}
Cauchy's momentum equation with constant density, constant viscosity, and no volume forces reduces
to the steady Navier--Stokes system \cite{bird1960transport},
\begin{equation}
\rho\,(\bu\cdot\nabla)\bu = -\nabla p + \mu\nabla^2\bu \qquad \text{in } \Om_f .
\label{eq:momentum}
\end{equation}
Assumptions A2 and A6 (constant Newtonian properties, laminar regime) make \eqref{eq:momentum}
well posed together with the boundary conditions of Chapter~\ref{ch:interface}. Every term carries
units \si{\pascal\per\metre}. Buoyancy would enter through the Boussinesq body force
$\rho\beta_f\theta_f\mathbf g$, where $\theta_f=T_f-T_{\mathrm{ref}}$ is the fluid temperature rise
and $\beta_f$ the volumetric expansion coefficient. The Richardson number computed in
Chapter~\ref{ch:baseline} on a length-consistent basis is $1.7\times10^{-3}$, so the Boussinesq
buoyancy term is excluded by the standard order-of-magnitude argument: the buoyant forcing is three
orders of magnitude below the inertial forcing. A Boussinesq body-force option is nevertheless kept
dormant in the formulation (one weak-form term) so that any future operating point with a
significantly larger fluid-side temperature excursion can be re-examined; for context on the
opposite limit, in which stable stratification measurably modifies channel heat transfer at bulk
Richardson numbers of order $10$ and above, see \cite{draskic2025stability}. The disposition is
re-examined in the limitations chapter.

\section{Energy conservation in the fluid}
For a constant-property fluid, the thermal energy equation reduces to the steady
advection--diffusion form \cite{bird1960transport}
\begin{equation}
\rho c_p\,(\bu\cdot\nabla)T_f = k_f\nabla^2 T_f \qquad \text{in } \Om_f .
\label{eq:fluidenergy}
\end{equation}
Viscous dissipation $\mu\Phi$, with $\Phi=2\mathbf D:\mathbf D$ and
$\mathbf D=\tfrac12(\nabla\bu+\nabla\bu^{\mathsf T})$ the rate-of-deformation tensor, is neglected
under Assumption~A7: at the baseline the Brinkman number is $1.0\times10^{-6}$
(Section~\ref{sec:baselinecalc}), and the dissipation term is required to be restored whenever a
computed case reaches $\mathrm{Br}\ge10^{-2}$. Each term has units \si{\watt\per\cubic\metre}.

\section{Heat conduction in the solid}
Steady conduction in the wall obeys \cite{carslaw1947conduction}
\begin{equation}
\nabla\cdot(k_s\nabla T_s)=0 \qquad \text{in } \Om_s^0 .
\label{eq:solidenergy}
\end{equation}
In the general theory of thermoelasticity the energy equation contains the deformation-heating
source $-\beta T_0\tr\dot{\beps}$, where $T_0$ is the absolute reference temperature
\cite{boley1960theory}. This source vanishes identically in any steady state
($\dot{\beps}\equiv0$), so \eqref{eq:solidenergy} is exact for the baseline. In a transient
start-up study its omission would be an approximation controlled by the thermoelastic coupling
number $\xi=\beta^2T_0/[(\lambda+2\mu_L)\rho_sc_s]$, evaluated at the baseline in
Section~\ref{sec:baselinecalc} ($\xi\approx2\times10^{-2}$).

The solid domain $\Om_s^0$ is the reference configuration: under small strains (A4) the distinction
between current and reference coordinates is negligible for the conduction problem. The wall thermal
field feeds the mechanical problem through the constitutive relation \eqref{eq:constitutive}, and
is closed by the conjugate interface conditions of Chapter~\ref{ch:interface}.

\section{Linear thermoelastic constitutive relations and kinematics}
The solid follows the Duhamel--Neumann extension of Hooke's law for an isotropic body with linear
thermal expansion \cite{boley1960theory,timoshenko1970elasticity},
\begin{equation}
\beps=\tfrac12\big(\nabla\bd+\nabla\bd^{\mathsf T}\big),\qquad
\bsig = 2\mu_L\beps + \lambda(\tr\beps)\bI - \beta\theta\bI
\qquad \text{in } \Om_s^0 .
\label{eq:constitutive}
\end{equation}
The assumptions are small strain (A4) and isotropy (A3). All stress components reported from the
model are \emph{total} Cauchy stresses, i.e.\ \eqref{eq:constitutive} is never split for reporting
into a ``mechanical'' and a ``thermal'' part; the trace of \eqref{eq:constitutive} is examined in
Section~\ref{sec:meanstress} and turns this convention into a checkable identity.

\subsection{Plane-strain closure and out-of-plane stress}
Under plane strain $\eps_{zz}=0$ (A1). Writing Hooke's law in compliance form with the free thermal
strain,
\begin{equation}
\eps_{zz}=\frac{1}{E}\big[\sigma_{zz}-\nu(\sigma_{xx}+\sigma_{yy})\big]+\alpha\theta=0
\;\;\Longrightarrow\;\;
\sigma_{zz}=\nu(\sigma_{xx}+\sigma_{yy})-E\alpha\theta ,
\label{eq:sigmazz}
\end{equation}
which specialises the general plane-strain result \cite{timoshenko1970elasticity,boley1960theory}
to the present notation. The equivalence of \eqref{eq:sigmazz} with the stiffness-form value
$\sigma_{zz}=\lambda\tr\beps-\beta\theta$ is verified algebraically in
Section~\ref{sec:deriv:planestrain}. The out-of-plane stress must be retained in any field
quantity that involves the stress trace, as the next section makes precise.

\section{Elastic equilibrium}
Neglecting inertia (A5) and volume forces, linear-momentum balance in the solid reduces to
\begin{equation}
\nabla\cdot\bsig=\mathbf{0} \qquad \text{in } \Om_s^0 .
\label{eq:equilibrium}
\end{equation}
Together with \eqref{eq:constitutive} and the stress boundary conditions of
Chapter~\ref{ch:interface}, \eqref{eq:equilibrium} constitutes the classical displacement
formulation of isothermal thermoelasticity \cite{boley1960theory}; elimination of stress gives the
Navier equation with the thermal body-force-equivalent term $-\beta\nabla\theta$.

\section{Mean stress and the spatial mean-stress identity}
\label{sec:meanstress}
The mean (hydrostatic) Cauchy stress $\sigma_m=\tr\bsig/3$ is used to characterise the compressive
membrane state of the wall. Two equivalent expressions follow from \eqref{eq:constitutive} and
\eqref{eq:sigmazz} and are derived in Section~\ref{sec:deriv:meanstress}:
\begin{equation}
\sigma_m=\frac{1+\nu}{3}\big(\sigma_{xx}+\sigma_{yy}\big)-\frac{E\alpha}{3}\,\theta
=K\big(\tr\beps-3\alpha\theta\big),
\label{eq:meanstress}
\end{equation}
where the second form identifies mean stress with bulk modulus $\times$ mechanical volumetric
strain. The two forms must agree identically at every point; this identity is used as an algebraic
consistency check in the post-processing of the future solver, and only complete thermoelastic
stress components may enter the first form. An immediate consequence for sanity checks
(Section~\ref{sec:deriv:sanity}) is that under plane strain even a \emph{free} uniformly heated
strip carries $\sigma_m=-E\alpha\theta/3\neq0$, because $\eps_{zz}=0$ is itself a constraint that
generates $\sigma_{zz}=-E\alpha\theta$; this corrects the three-dimension-free expectation
$\sigma_m=0$ recorded in an earlier draft of the formulation (Appendix~B, register D-2).

\section{Reduced wall equation for design estimates}
\label{sec:beam}
For design estimates and beam-level benchmarks, the strip is additionally approximated as a thin
Euler--Bernoulli plate strip of flexural rigidity per unit depth
$D_s=Ed^3/[12(1-\nu^2)]$ \cite{timoshenko1970elasticity}, governed by
\begin{equation}
D_s\,w^{''''}(x)= -\,p(x), \qquad
w(0)=w(L)=w^{''}(0)=w^{''}(L)=0,
\label{eq:beam}
\end{equation}
with $w$ measured positive into the channel (Section~\ref{sec:signs}) and $p(x)$ the fluid pressure
acting on the wall, pushing the wall \emph{away} from the channel; the sign convention is fixed so
that the pinned--pinned thermal bowing of Section~\ref{sec:deriv:deflection} is positive.
Equation \eqref{eq:beam} is used \emph{only} for design estimates and benchmark construction; the
production model is the two-dimensional system \eqref{eq:solidenergy}--\eqref{eq:equilibrium}, of
which \eqref{eq:beam} is the standard leading-order reduction for $d/L\ll1$
\cite{timoshenko1970elasticity}.

\section{The deformed fluid domain}
\label{sec:ale}
The fluid problem is posed on the domain $\Om_f$ occupied by the fluid in the converged steady
configuration. Given the solid solution on the reference interface, a mesh-displacement field
$\widehat{\mathbf w}$ is constructed on $\Om_f^0$ by solving a pseudo-elastic extension problem
(cf.\ Section~\ref{sec:coupling}) with Dirichlet data $\widehat{\mathbf w}=\bd$ on
$\Ga_{\mathrm{if}}^0$ and $\widehat{\mathbf w}=\mathbf 0$ on rigid boundaries, and
$\Om_f=\{X+\widehat{\mathbf w}(X):X\in\Om_f^0\}$. Because time does not appear in the baseline
problem, no mesh velocity enters the fluid equations and no geometric-conservation-law constraint
arises; such constraints become mandatory only for the transient extension, where the ALE
formulation must be made GCL-compliant \cite{lesoinne1996geometric}.

\section{Constitutive completeness}
No strain-gradient, nonlocal, fractional-derivative, memory, or two-temperature terms appear in
this model, by deliberate restriction (cf.\ Section~\ref{sec:limitations}). Any such addition would
require its own constitutive evidence, boundary conditions, and validation objects, none of which
are available in the literature base used here.
